\scanpage{7}
\uline{I bis}\qquad \uline{Exemples sur} $k=\mathbb{R}$

Prenons
\[
E=\mathbb{C} \tag*{(1)}
\]
regardé comme plan sur les réels. \deleted{\unsure{Soient}}Soit
\deleted{$\unsure{\zeta|z}$} $\zeta$ une racine \uline{primitive} \unsure{$n$:ième|$n$-ième} de 1 $\unsure{\big\}|:|,}$
\deleted{posons} (racine de 1 d'ordre $n$)
\[
\zeta^n=1,\qquad \zeta^d\ne 1\quad\text{si}\quad d\mid n,\deleted{\unsure{d|1}}1\leq d<n
\tag*{\deleted{\unsure{(a)|(2)}}(2)}
\]
posons
\[
s_i=\zeta^i\qquad i\in\mathbb{Z}/n\mathbb{Z}\tag*{(3)}
\]
avec la structure polygonale \inserted{combinatoire} standard.
\inserted{Les $s_i$ distincts car $\zeta$ racine \uline{primitive}.}
Soient $\sigma,u\in\deleted{\mathit{Aff}}\mathit{Gl}_{\mathbb{R}}(E)\subset\mathit{Aff}(E)$ définis par
\[
\left\{\begin{aligned}
\sigma z&=\bar z &&\unsure{\sigma|z}\in\\
uz&=\zeta z
\end{aligned}\right.
\tag*{\(\begin{aligned}&\text{(4)}\\&\deleted{\unsure{(5)}}\end{aligned}\)}
\]
donc on a
\[
\left\{\begin{aligned}
\sigma s_i&=s_{-i}\\
us_i&=s_{i+1}
\end{aligned}\right.\tag*{(5)}
\]
donc on a un $n$-polygone régulier, $P_{\zeta}$.

On vérifiera \deleted{que} aisément (et on retrouvera
ça dans un contexte plus général)
que
a) $P_\zeta$ isom \unsure{à|:} $P_{\zeta'}$ ssi $\zeta'=\zeta$ ou $\zeta'=\zeta^{-1}$
(ou encor, posant $\unsure{\alpha(\zeta)|\alpha()}=\zeta+\zeta^{-1}$, ssi $\alpha(\zeta)=\alpha(\zeta')$)

b) Tout polygône régulier \deleted{dans \illegible{1}}
un plan affine sur $\mathbb R$ est isom. à un $P_\zeta$.
% archive sheet number: 6

